%قبل از شروع مسئله ترجیحاً چیزی ننویسید. اما اگر خواستید بنویسید، در قالب لیستی (فقرات و شمارش) بنویسید.
%برای اضافه کردن نمره، بعدی از شروع‌مسئله، [x] را اضافه کنید و به جای x نمره مورد نظر را قرار دهید. برای زیرمسئله‌ها هم به همین صورت است.
%\شروع‌مسئله[10]
\شروع‌مسئله [15]
%\متن‌سیاه{سوال عملی اول:}\\
%\متن‌سیاه{دراپ‌تست}\\

%\متن‌سیاه{ورودی:}\\
%در خط اول ورودی، ابتدا k و سپس m داده می‌شود.\\
%\متن‌سیاه{خروجی:}\\
%در تنها سطر خروجی، تعداد آیفون‌های مورد نیاز می‌آید. اگر در هیچ صورتی چنین کاری ممکن نبود ۱- چاپ می‌شود.\\
%\متن‌سیاه{محدودیت‌ها:}\\
%$1\leq k\leq 10^4$

%$m\leq10^9$
%\\
%\textbf{\متن‌سیاه{ورودی و خروجی نمونه}\\
%% تست اول
%\begin{latin}
% \begin{tabularx}{\textwidth}{|X | X|}
% \hline
% \textbf{stdin} &
% \textbf{stdout}\\
% \hline
% 100 100&
% 1\\
% \hline
% 
% \end{tabularx}
%\end{latin}
%
%% تست دوم
%
%\begin{latin}
% \begin{tabularx}{\textwidth}{|X | X|}
% \hline
% \textbf{stdin} &
% \textbf{stdout}\\
% \hline
% 340 20&
% 3\\
% \hline
% 
% \end{tabularx}
%\end{latin}}
% صورت اصلی مسئله را در این‌جا بنویسید.

فرض کنید $U=\{ \, (x_1,x_2,\dots ,x_n) \mid x_i \in \{0,1\}\}$ و $V=\{ \, (x_1,x_2,\dots ,x_m) \mid x_i \in \{0,1\}\}$ دو فضای برداری روی میدان ${\mathbb{Z}}_2$ باشند. جمع برداری و ضرب اسکالر در این میدان همان جمع و ضرب عادی است که به پیمانۀ 2 انجام می‌شود.

اگر $m<<n$، می‌توانیم تبدیل خطی $T:U\rightarrow V$ را یک تابع درهم‌سازی درنظر بگیریم که هر عضو $U$ را به عضوی از $V$ می‌برد.
$\mathcal{H}$ را خانوادۀ تمام تبدیل‌های خطی از $U$ به $V$ بگیرید.
ثابت کنید $\mathcal{H}$ سراسری \متن‌لاتین{(universal)} است.

به عبارت دیگر باید ثابت کنید که اگر $h$ یک عضو $\mathcal{H}$ باشد که به صورت تصادفی با توزیع همگن انتخاب شده است،
$$\Pr ( h(x) = h(y) ) \le \frac{1}{|V|} \quad  \qquad \forall x\neq y \in U$$
که $|V| = 2^{m}$.

\solution{
هر تابع خطی در
$\mathcal{H}$
را می‌توان به صورت یک ماتریس
$m \times n$
با درایه‌های صفر و یک نمایش داد. دو بردار
$x = ( x_1, x_2, \dots , x_n)$
و
$y = (y_1, y_2, \dots , y_n)$
در $U$ را در نظر بگیرید که $x \neq y$. می‌خواهیم تعداد ماتریس‌های $T$ با درایه‌های صفر و یک $m \times n$ را بشماریم که
$T(x) = T(y)$.

$x \neq y$.
در نتیجه $k$ای وجود دارد که
$x_k \neq y_k$.
ادعا می‌کنیم که در هر حالتی که تمام ستون‌های ماتریس $T$ به غیر از ستون $k$ام، می‌توانند داشته باشند، ستون $k$ام دقیقا یک حالت می‌تواند داشته باشد تا تساوی $T(x) = T(y)$ برقرار شود. فرض کنید
$t_{ij}$
درایه سطر $i$ام و ستون $j$ام $T$ باشد. همچنین فرض کنید
$T(x) = T(y)$.
برای هر $i$، $1 \le i \le n$ داریم:
$$\sum_{j = 1}^{n} t_{ij}x_j = \sum_{j = 1}^{n} t_{ij}y_j$$
بنابراین درایه $i$ام ستون $k$ام برابر است با:
$$t_{ik} = (x_k - y_k)^{-1} \sum_{j = 1, j \neq k}^{n} t_{ij}(y_j - x_j)$$
ستون $k$ام به طور یکتا مشخص شد. پس از بین تمام
$2^{mn}$
تابع خطی در
$\mathcal{H}$،
دقیقا
$2^{mn}/2^m = 2^{mn - m}$
تابع هستند که
$T(x) = T(y)$.
اکنون:
$$\Pr ( h(x) = h(y) ) = \frac{2^{mn - m}}{2^{mn}} = \frac{1}{2^m} = \frac{1}{|V|}$$
و حکم مورد نظر ثابت می‌شود.
	
}
%%%%%%%%%%%%%%%%%%%%%%%%%%%%%%%%%%%%%%%%%%%%%%%%%%%%%
% برخی دستورات مفید برای تمرین‌ها. %
%%%%%%%%%%%%%%%%%%%%%%%%%%%%%%%%%%%%%%%%%%%%%%%%%%%%%

% فرمول‌های نمونه:
%$f(x) = \Theta(x)$
%$T(x, y) \leq \log(x) + 2^y + \binom{2n}{n}$
% {\tt code\_function}


% لیست موارد بدون شماره‌گذاری:
%\شروع{فقرات}
% \فقره مورد اول 
% \شروع{فقرات}
% \فقره مورد اول داخلی 
% \شروع{فقرات}
% \فقره اولین مورد خیلی داخلی 
% \فقره دومین مورد خیلی داخلی 
% \پایان{فقرات}
% \فقره مورد دوم داخلی 
% \پایان{فقرات}
% \فقره مورد دوم 
%\پایان{فقرات}

% لیست موارد شماره‌گذاری شده:
%\شروع{شمارش}
% \فقره یک مورد 
% \فقره یک مورد دیگر 
% \فقره باز هم یک مورد دیگر
%\پایان{شمارش}

% برای نوشتن یک فرمول چند خطی که یک شماره بخوره. برای حالت‌های دیگر به سایت زیر مراجعه کنید:
% https://www.overleaf.com/learn/latex/Aligning_equations_with_amsmath
%\begin{align} 
% \begin{split}
% (x+y)^3 &= (x+y)^2(x+y) \\
% &= (x^2+2xy+y^2)(x+y) \\
% &= (x^3+2x^2y+xy^2) + (x^2y+2xy^2+y^3) \\
% &= x^3+3x^2y+3xy^2+y^3
% \end{split} 
%\end{align}

% نوشتن ماتریس و آرایه دوبعدی
%\begin{align}
% A = 
% \begin{bmatrix}
% A_{11} & A_{21} \\
% A_{21} & A_{22}
% \end{bmatrix}
%\end{align}

% اضافه کردن تصویر با شرح. برای اضافه کردن تصویر بدون شرح، فقط دستور درج‌تصویر را استفاده کنید.
%\شروع{شکل}[h!]
%\تنظیم‌ازوسط
% \درج‌تصویر[پهنا=0.5\پهنای‌متن]{img.jpg}
%\شرح{شرح تصویر}
%\برچسب{تصویر:۱}
%\پایان{شکل}
%\پایان{زیرمسائل}

%کد زیر را ببینید.
%%\lstinputlisting{source_filename.py}
%\begin{latin}
%\begin{lstlisting}
%# Import the modules
%import sys
%import random
%
%ans = True
%
%while ans:
%question = raw_input("Ask the magic 8 ball a question: (press enter to quit) ")
%
%answers = random.randint(1,8)
%
%if question == "":
%sys.exit()
%
%elif answers == 1:
%print "It is certain"
%
%elif answers == 2:
%print "Outlook good"
%
%elif answers == 3:
%print "You may rely on it"
%
%elif answers == 4:
%print "Ask again later"
%
%elif answers == 5:
%print "Concentrate and ask again"
%
%elif answers == 6:
%print "Reply hazy, try again"
%
%elif answers == 7:
%print "My reply is no"
%
%elif answers == 8:
%print "My sources say no"
%\end{lstlisting}
%\end{latin}